% SPDX-License-Identifier: LPPL-1.3c
% Copyright (C) 2026 Christophe Jorssen
% Author and Current Maintainer: Christophe Jorssen <christophe.jorssen@gmail.com>
% LPPL maintenance status: maintained
% This file is part of luafemm. See LICENSE and LICENSES.md for its terms.
\section{Solving, field lines, mesh and numeric results}
\label{sec:field}
\begin{key}{/tikz/femm/solve}
Resolve the current problem when this path executes, for example
|\path[femm/solve];|. This is optional: field pictures and requested
profile values solve on demand. Declare all regions before the first
mesh or solve request. A mesh-only picture does not solve the system.
Section~\ref{sec:workflow} gives the complete declaration and computation order.
\end{key}
\begin{command}{\femmsolve}
Command form of the same solve operation. A current model must exist.
\end{command}

\begin{key}{/tikz/pics/femm field}
Used as |\pic {femm field};|. Solve if necessary and draw oriented
isolines of $A_z$. Each level can have several connected components.
Because $\mathbf B=(\partial_y A_z,-\partial_x A_z)$, these are field lines
of the computed P1 solution. They are polygonal, without graphical smoothing.
\end{key}
\begin{key}{/tikz/pics/femm mesh}
Used as |\pic {femm mesh};|. Generate the mesh if necessary, then draw each
edge once. It is useful for inspecting interfaces and local refinement.
\end{key}
Both pictures are anchored to the original model frame. Moving the pic
inside a local scope does not detach the field from the media. Use a clip
to display a smaller part of the computed domain.

\begin{key}{/tikz/femm/lines=\meta{positive integer} (initially 29)}
Number of evenly spaced potential levels. Levels are placed at the
midpoints of equal bins between the minimum and maximum potential.
This counts levels, not necessarily visible connected curves. Zero-field
models produce no contours.
\end{key}
\begin{key}{/tikz/femm/arrows=\meta{true or false} (initially true)}
Enable or disable the direction markers. The direction follows B, not an
arbitrary contour traversal. With no value, PGF's boolean handler selects true.
\end{key}
\begin{key}{/tikz/femm/arrow spacing=\meta{dimension} (initially 22mm)}
Distance between repeated markers along the rendered path. Markers occupy
the middle part of each contour (from relative positions .12 to .95).
Short paths need not contain a marker. This is a graphic dimension, not
a mesh length or physical parameter.
\end{key}
\begin{key}{/tikz/femm/arrow length=\meta{dimension} (initially 1.3mm)}
Length of the Stealth marker, in the displayed figure.
\end{key}
\begin{key}{/tikz/femm/arrow width=\meta{dimension} (initially .85mm)}
Width of the Stealth marker.
\end{key}
\begin{key}{/tikz/femm/field style=\marg{TikZ options}}
Extra graphics options for this field picture; empty initially. For
example |{teal!70!black,line width=.5pt}|. The base |femm field| style is
applied first, and these options second. This key sets rather than appends
the extra options within the current TeX scope.
\end{key}
\begin{key}{/tikz/femm/mesh style=\marg{TikZ options}}
Extra mesh graphics options, empty initially. The base |femm mesh| style
is applied first. No material or mesh-size property is changed by a
graphic style.
\end{key}
\begin{stylekey}{/tikz/femm field}
Base contour style: |blue!65!black,line width=.32pt,femm oriented field|.
Extend with |/.append style| if the standard arrow handling is to remain.
\end{stylekey}
\begin{stylekey}{/tikz/femm mesh}
Base mesh style: |draw=black!25,line width=.12pt|.
\end{stylekey}
\begin{key}{/tikz/femm oriented field}
Reusable marker setup used by |femm field|. It respects |femm/arrows|
and the marker dimensions. PGF's floating-point vector length is enabled
locally for decoration: very short contour segments would otherwise
overflow PGF's fixed-point reciprocal. The contour vertices are unchanged.
\end{key}
The helper styles |femm current field| and |femm current mesh| are internal
storage for the corresponding extra-style keys; use those keys rather
than depending on the helper implementation.

\subsection{Numeric output}
\begin{command}{\femmB\marg{x}\marg{y}}
Print $|\mathbf B|$ in tesla, with three digits after the decimal point.
Coordinates are model-unit numeric literals. The current model must
already be solved; the first containing triangle is used at interfaces.
For more precision or components, use a profile or the Lua API.
\end{command}
\begin{command}{\femmstat\marg{name}}
Print a solved-model statistic with up to eight significant digits.
Unknown or unavailable values are errors.
\end{command}
\begin{center}
\begin{tabular}{ll}
\toprule Name & Meaning \\
\midrule
|nodes|, |triangles| & Mesh vertex and element counts \\
|newton| & Accepted Newton corrections \\
|cg| & Total conjugate-gradient iterations \\
|residual| & Final relative nonlinear residual \\
|seconds| & CPU solution time, excluding mesh generation \\
|peak| & Largest element value of $|\mathbf B|$ in T \\
|net_current| & Integral of $J_z$ over the section, in A \\
\bottomrule
\end{tabular}
\end{center}
\begin{command}{\femmmeshstat\marg{name}}
Print a constrained-mesh diagnostic with up to six significant digits.
These values require the Delaunay mesher; they do not exist for |grid|.
\end{command}
\begin{center}
\begin{tabular}{ll}
\toprule Name & Meaning \\
\midrule
|min_angle| & Smallest triangle angle in degrees \\
|max_area| & Largest triangle area in model units squared \\
|mesh_tolerance| & Actual vertex-merging distance in model units \\
|seconds| & CPU mesh generation and audit time \\
|flips| & Constraint and refinement legalisation flips \\
|steiner_points| & Additional quality-refinement vertices \\
|segment_splits| & Quality-refinement segment bisections \\
|circumcenters| & Accepted quality-refinement circumcentres \\
|exact_orientation|, |exact_incircle| & Exact-arithmetic fallback counts \\
|exact_diametral| & Exact diametral-disk fallback count \\
\bottomrule
\end{tabular}
\end{center}
Statistics remain accessible after the source picture's TeX group closes.
Starting another model changes the current model for these commands;
named profiles, in contrast, retain their own source model.
